\chapter{Installation}
\label{ch:installation}

\jns{} is a pure-Python package that drives external programs; it bundles
none of them.  The core installation needs only \file{ase} and \file{numpy}
and already covers the calculators, the uncertainty-quantification
mathematics and the calibration tools.  The deep-learning stack and the
SCINE engine integration are optional extras, so an installation that only
needs an ORCA or xTB calculator does not have to carry them.

\section{Requirements}

\begin{table}[htbp]
  \centering
  \small
  \begin{tabularx}{\linewidth}{@{}>{\raggedright\arraybackslash}p{3.4cm}
    >{\raggedright\arraybackslash}p{3.9cm} X@{}}
    \toprule
    Component & Form & Needed for \\
    \midrule
    Python & $\ge$ 3.11 & everything \\
    \file{ase} $\ge$ 3.22, \file{numpy} $\ge$ 1.24 & installed with the core & all calculators, the UQ module \\
    \file{xtb} & external executable, optional & L1 pre-screening (\code{XTBCalculator}) \\
    ORCA & external program, version 6.x, academic license & L2/L3, fallback of the safe calculator \\
    \file{torch}, \file{mace-torch}, \file{aimnet2} & \code{[nnp]} extra & NNP committee, \code{SafeNNPCalculator} \\
    \file{scine-utilities}, \file{scine-chemoton}, \file{scine-database}, \file{scine-puffin} & \code{[scine]} extra & Chemoton/Puffin integration \\
    MongoDB & running server & Chemoton database layer \\
    \file{pytest} $\ge$ 7 & \code{[dev]} extra & running the test suite \\
    \bottomrule
  \end{tabularx}
  \caption{Requirements by usage.  The first two rows are sufficient for the
  calibration workflow of path~1 in Chapter~\ref{ch:getting-started}.}
\end{table}

\section{Installing the package}

\begin{codeblock}
pip install jnuscout                  # core: calculators (ORCA/xtb), UQ maths
pip install "jnuscout[nnp]"           # + torch, mace-torch, aimnet2
pip install "jnuscout[scine]"         # + scine-utilities/chemoton/puffin/database
pip install "jnuscout[dev]"           # + pytest
\end{codeblock}

Extras can be combined, for example \code{pip install "jnuscout[nnp,scine]"}.
Quoting matters in shells that expand brackets (\file{zsh} and some
\file{bash} configurations); the quotes above are not optional there.

\begin{table}[htbp]
  \centering
  \small
  \begin{tabularx}{\linewidth}{@{}l X@{}}
    \toprule
    Extra & Adds \\
    \midrule
    (core) & \file{ase} $\ge$ 3.22, \file{numpy} $\ge$ 1.24 \\
    \code{[nnp]} & \file{torch} $\ge$ 2.0, \file{mace-torch}, \file{aimnet2} \\
    \code{[scine]} & \file{scine-utilities}, \file{scine-chemoton}, \file{scine-database}, \file{scine-puffin} \\
    \code{[dev]} & \file{pytest} $\ge$ 7 \\
    \bottomrule
  \end{tabularx}
  \caption{Dependency extras defined in \file{pyproject.toml}.}
\end{table}

\section{External programs}

\subsection{ORCA (L2, L3 and the fallback)}

Levels L2 and L3, and every fallback of the safe calculator, are single
ORCA runs.  \jns{} neither bundles nor downloads ORCA: obtain version 6.x
from the developers under an academic license and make the executable
available either on \code{PATH} or through the \code{ORCA_BIN} environment
variable (default: \file{orca}).  The constructor keyword \code{orca_bin=}
overrides both for one calculator instance.

Two resource settings are forwarded to every ORCA input: \code{nprocs} and
\code{maxcore} (memory in MB per core), written as the directives
\%\code{pal nprocs} and \%\code{maxcore}.  Their defaults are 8 and 4000.

ORCA is started in an adjusted process environment.  The calculator prepends
the ORCA directory, the directories listed in \code{ORCA_EXTRA_LIB_DIRS}
(colon-separated) and the common OpenMPI library locations
(\file{/usr/local/openmpi-4.1.8/lib}, \file{/usr/lib64/openmpi/lib},
\file{/usr/lib/x86_64-linux-gnu/openmpi/lib}) to \code{LD_LIBRARY_PATH},
and the corresponding \file{bin} directories plus
\code{ORCA_EXTRA_BIN_DIRS} to \code{PATH} so that \file{mpirun} can be
resolved.

\begin{notebox}[ORCA exits without writing a file]
If ORCA terminates with return code 0 but produces no \file{.engrad} file,
the cause is usually the launcher environment rather than the calculation.
A non-interactive shell (cron, \file{systemd-run}) does not read
\file{.bashrc}, so the library path a login shell would have set is missing
and the binary cannot load \file{liborca} or \file{libmpi}.  If your OpenMPI
installation lives elsewhere, set \code{ORCA_EXTRA_LIB_DIRS} and
\code{ORCA_EXTRA_BIN_DIRS} to its library and binary directories.
\end{notebox}

\subsection{xTB (L1, optional)}

L1 pre-screening uses the \code{xtb} command-line program and is optional:
the package works without it as long as no L1 calculation is requested.
Point \code{XTB_BIN} (default \file{xtb}) at the executable, or pass
\code{xtb_bin=} to the calculator.  Any build exposing the standard command
line (\code{--gfn}, \code{--chrg}, \code{--uhf}, \code{--grad}) is usable;
the conda-forge build compiled with gfortran is the one tested.  The wrapper
exists because ASE 3.29 ships no built-in xTB interface and the SCINE xTB
module was not usable on the deployment host (it requires an AVX-512 capable
CPU, and its ABI conflicted with the installed SCINE release).

\section{Model weights and offline operation}

The committee (Chapter~\ref{ch:calculators}) has three members, loaded
lazily on first use.  Their identifiers are read from the environment:

\begin{table}[htbp]
  \centering
  \small
  \begin{tabularx}{\linewidth}{@{}l l X@{}}
    \toprule
    Member & Variable & Default \\
    \midrule
    MACE-MP-0 & \code{MACE_MP0_PATH} & \file{medium} (resolved and, if absent, downloaded by \file{mace-torch}) \\
    MACE-MPA-0 & \code{MACE_MPA0_PATH} & \url{~/.cache/mace/mace-mpa-0-medium.model} \\
    AIMNet2-b973c & --- & selected by model identifier \file{aimnet2_b973c}; no path variable in this version \\
    \bottomrule
  \end{tabularx}
  \caption{Weight resolution for the default committee members.}
\end{table}

\begin{notebox}[Offline machines]
If the MACE weights are not cached locally, \file{mace-torch} attempts a
download on first use, which fails (or stalls) without network access.  Point
\code{MACE_MP0_PATH} and \code{MACE_MPA0_PATH} at local \file{.model}
files instead; any path is accepted, not only the cache directory.
Everything else in \jns{} runs locally: no structure ever leaves the machine.
\end{notebox}

\begin{notebox}[When environment variables are read]
\code{ORCA_BIN} and \code{XTB_BIN} are read when
\file{jnuscout.calculators.orca} and \file{jnuscout.calculators.xtb} are
imported, and the MACE paths when
\file{jnuscout.calculators.committee} is imported.  Set the variables before
starting the interpreter, or use the per-calculator keywords
(\code{orca_bin=}, \code{xtb_bin=}, \code{members=}), which take effect at
any time.
\end{notebox}

\section{The SCINE engine and MongoDB}

The \code{[scine]} extra installs the engine stack.  Chemoton keeps its
reaction network in a MongoDB database, and the stock behaviour is to expect
a running \file{mongod}; \jns{} does not change this.  The bridge and the
injection patches were developed against Chemoton 4.1.0 with
\file{scine_utilities} 10.1.0; other releases of the SCINE stack are not
covered by the current test suite.  How the NNP is registered as a method
family is described in Chapter~7.

\section{Installing from source}

\begin{codeblock}
git clone https://github.com/jqChen1566/jnu-scout.git
cd jnu-scout
pip install -e ".[nnp]"
bash scripts/verify.sh
\end{codeblock}

\file{scripts/verify.sh} runs the unified check used for releases: a syntax
gate (\code{python -m compileall} over \file{src}, \file{scripts} and
\file{tests}) followed by the test suite.  A source checkout has no compiled
components; the package is Python only.

\section{Checking the installation}

The two commands below exercise the package without launching any external
program (a calculator runs its backend only inside a calculation):

\begin{transcript}
$ python -c "import jnuscout; print(jnuscout.__version__)"
0.1.0
$ python -c "from jnuscout.calculators import CalculatorFactory; print(CalculatorFactory.create('L1').method)"
GFN2-xTB
\end{transcript}

If the first command fails, the package is not on the import path of the
interpreter you are using.  If the second fails while the first succeeds, the
core dependencies are broken; reinstall the package.
